\newpage
\subsubsection{Full-Search research strategy}
We consider windows of the image, with size $=(2A+1)\times (2B+1)$.\\
Then we build a $(2A+1)\times (2B+1)$ SSD matrix, this brings $\propto n^2$ complexity.
By searching the minimum of the SSD matrix, it is the best possible vector that can be used to approximate the window.
\subsubsection{Fast Research Methods: 3SS (Three step search, 2D-log)}
The use of Full-search is computationally expensive, we need strategies to test less positions in the window. Assuming error function is unimodal $\Rightarrow \exists$ one global minimum.\\
3SS method allows to test only 9 positions ($x \in \{0, \pm A/2\}$,$y \in \{0, \pm B/2\}$).Avoiding to compute $89\%$ operations.\\
We recursively apply this method on the sub-windows centered in the minimum valued element of the matrix, and also reducing the size of the sub-matrix by a half, until we get a 1-pixel subwindow.
\subsubsection{Fast Research Methods: Diamond Search}
We rotate the initial 9-points structure by 45 degrees, thus we use a diamond pattern to search, new iteration can cost from $3$ (diagonal displacement) to $5$ (horizontal/vertical displacements) positions to compute.\\
Recursively refining until final refinement is done through a cross-shaped 5-position small diamond.\\
Avoiding to compute $90\%$ operations.
\subsubsection{Fast Research Methods: Hex Search}
As the diamond search, here we use an hexagon-like shape with 7 positions to compute. Horizontal, vertical and diagonal displacements cost 3 positions to compute.
Final refinement is done through the same cross-like shape of the diamond search. Avoiding to compute $92\%$ operations.
\subsubsection{Fast Research Methods: TZSearch}
Robust to local minima, phases:
\begin{itemize}
    \item Search predictors: search time/space neighbors, if low error stop
    \item Adaptive loop: if no predictor is good, use diamond/square search increasing steps
\end{itemize}
\subsubsection{BM Improvement: Sub-pixel precision}
With half pixel precision, we can reduce cost function because we have more options.\\
Problem: we have to increase resolution, we can use some interpolation methods, steps:
\begin{enumerate}
    \item select a certain pixel $(\hat{i}, \hat{j})$
    \item we test $(\hat{i}\pm \frac{1}{2}, \hat{j}\pm \frac{1}{2})$
    \item recursively we test all subpixels with $1/4, 1/8$...
\end{enumerate}
\subsubsection{BM Improvement: Variable Blocksize}
Foreach block $b$ of size $B$ in image $F$:
\begin{itemize}
    \item calculate $J = D + \lambda R$
    \item divide b into 4 subblocks
    \item compute $J_i$ $\forall i$ subblocks, $J_{sub} = \sum J_i$
    \item if $J_{sub} < J(v)$ apply algorithm to subblocks, else keep and store MV
\end{itemize}
\subsection{Parametric Methods}
We want MVF as a function of each pixel, degrees of freedom = parameters of the function.\\
\subsubsection{Affine Model}
$$\vec{v}(p) = \vec{b} + Bp = \left[\begin{matrix}
    b_1\\
    b_2
\end{matrix}\right] + \left[\begin{matrix}
    b_3 & b_4\\
    b_5 & b_6
\end{matrix}\right]p$$
Special case if $B=0 \to$ translation (Block Matching), other possible intermediate cases are:
\begin{itemize}
    \item Zoom In or Zoom Out
    \item Rotations
    \item Displacement (Translations)
\end{itemize}
\subsubsection{Best possible parameters}
Which are the best possible parameters? We know the affine model has 6 parameters $\pi = \left[b_1 \cdots b_6\right]$\\
\textbf{Indirect Estimation}
$$\pi^\star = \arg\min_{\pi}\sr{}{n,m\in R}\left[u(n,m)-u_\pi(n,m)\right]^2 + \left[v(n,m)-v_\pi(n,m)\right]^2$$
\textbf{Direct Estimation}, first approach
$$\pi^\star = \arg\min_{\pi}\sr{}{n,m\in R}\left[u_\pi(n,m)f_x(n,m) + v_\pi(n,m)f_y(n,m) + f_t(n,m)\right]^2$$
\textbf{Direct Estimation}, second approach (SAD/SSD minimization)
$$\pi^\star = \arg\min_{\pi}\sr{}{n,m\in R}\left[f(n-u_\pi(n,m),m-v_\pi(n,m), t-1) - f(n,m,t)\right]^2$$
\subsection{Deep-Learning for Motion Estimation}
\begin{itemize}
    \item Specialized architectures (CNNs)
    \item Large-scale datasets for training (Synthetic or Real, used by some game video quality enhancement tools)
    \item High computational power of inference by using GPUs (in the last few years)
\end{itemize}
\subsubsection{Time evolution}
\begin{itemize}
    \item First Generation: \textbf{FlowNet (2015)}, first CNN for motion estimation (input images, output optical flow)
    \item Second Generation: \textbf{FlowNet 2.0 (2017)}, better tradeoff complexity-performance
    \item Third Generation: \textbf{PWC-Net (2018)}, faster and quite efficient in complexity, easier to implement
    \item Nowadays:\textbf{RAFT (2020)}, slower than PWC-Net, precise on thin structures, iterative updates and GRU\footnote{Gated Recurrent Unit}.
\end{itemize}
\subsubsection{Usage Comparison}
DL methods has the best efficiency in analysis tasks (Optical Flow).\\
DL methods used in real-time video coding need to be tuned by practical constraints.\\
\textbf{Advantages}: Robustness, high precision\\
\textbf{Challenges}: Generalization on different datasets, computational cost to be optimized.
\section{Video-coding Principles}
We want to exploit both spatial redundancy and temporal redundancy.
\subsection{Block Schema}
$$x \longrightarrow \boxed{\text{Temporal Compression}} \longrightarrow \boxed{\text{Spatial Compression}} \longrightarrow \boxed{\text{Buffer}} \longrightarrow$$
Can we swap Temporal and Spatial compression blocks? Yes but it isn't used because it doesn't make sense.
\subsection{Block-Matching motion estimation}
$$J(v) = d(B_k^{(p)},B_h^{(p+v)}) + \lambda_{ME}R(v)$$
How can we build the prediction using motion estimation? Motion Compensation
\subsubsection{Motion Compensation}
$$\hat{I}_k(p) = I_h(p + v^\star(p))$$
Sometimes MC works better and sometimes it works worse, solution: adaptive blocks.
\subsubsection{Adaptive block coding}
\begin{itemize}
    \item Intra Coding: encode block as it is
    \item Inter Coding: use other images to predict it
\end{itemize}
$\forall B_k^{(p)}$ block of image k:
\begin{itemize}
    \item compute $J(v)$
    \item find the $v^\star = \arg\min_pv(p)$
    \item encode block:
    \begin{itemize}
        \item If intra: encode in a JPEG-like way
        \item If Inter:
        \begin{itemize}
            \item Decode $v^\star$ motion vector
            \item Decode predictor error blocks $E_p^{(b)}$
            \item Correct blocks
        \end{itemize}
    \end{itemize}
\end{itemize}
